\subsection{特征标公式}

在本号中，将群 \(X(T)\) 写成乘法记号，并将其元素视为 T 上的复值函数。假设元素 \(\rho\) 是 \(X(T) \otimes Q\) 中属于 \(X(T)\) 的元素。

将 \(L^{2}(T)\) 记为 T 上平方可积复函数的等价类组成的 Hilbert 空间，将 \(\Theta(T)\) 记为连续代表函数组成的子空间。根据《谱理论》，\(X(T)\) 是 \(L^{2}(T)\) 的正交归一基，也是 \(\Theta(T)\) 的代数基。

对于 \(f \in \mathrm{L}^{2}(\mathrm{T})\) 和 \(w \in W\)，将 \(\mathrm{L}^{2}(\mathrm{T})\) 中由 \(wf(t) = f(w^{-1}(t))\) 定义的元素 wf 记为；于是，对于 \(\lambda \in \mathrm{X}(\mathrm{T})\)，有 \(w\lambda = w(\lambda)\)。将 \(\varepsilon : W \to \{1, -1\}\) 记为符号（即满足对每个反射 s 都有 \(\varepsilon(s) = -1\) 的唯一同态）；对于 \(f \in \mathrm{L}^{2}(\mathrm{T})\)，置

\[
\mathrm{J} (f) = \sum_ {w \in \mathrm{W}} \varepsilon (w) ^ {w} f.
\]

若 \(\lambda \in X_{++}\)，则各特征标 \(^w (\lambda \rho)\) 彼此不同；事实上，只需证明 \(^w (\lambda \rho)\neq \lambda \rho\) 对所有 \(w\neq 1\) 成立；但这由引理 1（第 1 号）以及事实 \(\langle \lambda \rho ,K_{\alpha}\rangle = \langle \lambda ,K_{\alpha}\rangle +1 > 0\) 对每个正根 \(\alpha\) 成立得到。因此，

\[
\| \mathrm{J} (\lambda \rho) \| ^ {2} = \operatorname{Card} (\mathrm{W}) = w (\mathrm{G}).
\]

称 \(f \in \mathrm{L}^{2}(\mathrm{T})\) 为反不变元，若 \(^{w}f = \varepsilon(w)f\) 对所有 \(w \in W\) 成立（即，对每个反射 s 都有 \(^{s}f = -f\)）。我们将证明 \(\frac{1}{w(\mathrm{G})}\mathrm{J}\) 是 \(\mathrm{L}^{2}(\mathrm{T})\) 到反不变元子空间上的正交投影。事实上，令 \(f, f'\) 属于 \(\mathrm{L}^{2}(\mathrm{T})\)，且 \(f'\) 反不变；则 \(\mathrm{J}(f)\) 反不变，并且

\[
\begin{array}{c} \langle f ^ {\prime}, \mathrm{J} (f) \rangle = \sum_ {w \in \mathrm{W}} \varepsilon (w) \langle f ^ {\prime},   ^ {w} f \rangle = \sum_ {w \in \mathrm{W}} \langle^ {w} f ^ {\prime},   ^ {w} f \rangle \\ = \sum_ {w \in \mathrm{W}} \langle f ^ {\prime},   f \rangle = w (\mathrm{G}) \langle f ^ {\prime},   f \rangle . \end{array}
\]

命题 3. 元素 \(\mathrm{J}(\lambda\rho)/\sqrt{w(\mathrm{G})}\)，其中 \(\lambda\in X_{++}\)，构成 \(\mathrm{L}^{2}(\mathrm{T})\) 的反不变元子空间的正交归一基，并构成 \(\Theta(\mathrm{T})\) 的反不变元子空间的代数基。

证明与第 VI 章第 3 节第 3 号命题 1 的证明相同。

根据第 VI 章同上引文的命题 2，

\[
\mathrm{J} (\rho) = \rho \prod_ {\alpha > 0} (1 - \alpha^ {- 1}) = \rho^ {- 1} \prod_ {\alpha > 0} (\alpha - 1),\tag{3}
\]

因此

\[
\mathrm{J} (\rho) \overline {{\mathrm{J} (\rho)}} = \prod_ {\alpha} (\alpha - 1).\tag{4}
\]

由定理 1 的推论 2（第 6 节，第 2 号），得到：

引理 4. 若 \(\varphi\) 和 \(\psi\) 是 \(\mathbf{G}\) 上的两个连续中心函数，

\[
\int_ {\mathrm{G}} \overline {{\varphi (g)}} \psi (g) d g = \frac {1}{w (\mathrm{G})} \int_ {\mathrm{T}} \overline {{(\varphi (t) \mathrm{J} (\rho) (t))}}. (\psi (t) \mathrm{J} (\rho) (t)) d t.
\]

对于所有 \(\lambda\in X_{++}\)，将 \(\chi_{\lambda}\) 记为最高权为 \(\lambda\) 的 G 的不可约表示的特征标。

定理 2（H. Weyl）。对于所有 \(\lambda \in \mathrm{X}_{++}\)，有 \(\mathrm{J}(\rho)\chi_{\lambda}|\mathrm{T} = \mathrm{J}(\lambda \rho)\)。

函数 \(\mathrm{J}(\rho)\cdot\chi_{\lambda}|T\) 在 W 下反不变，并且是 \(\mathrm{X(T)}\) 中元素的整数系数线性组合。因此，根据第 VI 章第 3 节第 3 号命题 1，它可以写成 \(\sum_{\mu}a_{\mu}\mathrm{J}(\mu\rho)\)，其中 \(\mu\) 属于 \(X_{++}\)，且 \(a_{\mu}\) 是整数，除有限多个外均为零；由于 \(\int_{\mathrm{G}}|\chi_{\lambda}(g)|^{2}dg=1\)

（《谱理论》），由命题 3 和引理 4 得到 \(\sum_{\mu}(a_{\mu})^{2} = 1\)；因此 \(a_{\mu}\) 除一个外全为零，而该值必为 1 或 \(-1\)。但是，\(\lambda\) 在 \(\chi_{\alpha}|T\) 中的系数等于 1（定理 1），所以 \(\lambda \rho\) 在 \(J(\rho)\cdot \chi_{\lambda}|T\) 中的系数等于 1（第 VI 章，第 3 节，第 3 号，备注 2），这意味着 \(a_{\lambda \rho} = 1\)，定理得证。

推论 1. 按第 3 号的记号，在 \(\mathbf{Z}[\mathrm{X(T)}]\) 中有

\[
\left(\sum_ {w \in \mathrm{W}} \varepsilon (w) e ^ {w \rho}\right) \operatorname{Ch} [ \lambda ] = \sum_ {w \in \mathrm{W}} \varepsilon (w) e ^ {w \lambda} e ^ {w \rho} \text {for all} \lambda \in \mathrm{X} _ {+ +}.
\]

这由定理和交换图（2）（第 3 号）得到。

推论 2. 对所有 \(\lambda \in X_{++}\) 以及 T 的每个正则元 \(t\)，

\[
\chi_ {\lambda} (t) = \frac {\sum_ {w} \varepsilon (w) \lambda (w t) \rho (w t)}{\sum_ {w} \varepsilon (w) \rho (w t)}\tag{5}
\]

其中两个求和遍历 W 的元素 w。

事实上，\(\mathrm{J}(\rho)(t)\) 非零，因为 \(t\) 是正则元（公式（4））。

若 \(\varphi\) 是 \(\mathbf{G}\) 上的中心函数，则 \(\varphi\) 在 \(\mathrm{T}\) 上的限制在 \(\mathrm{W}\) 下不变，所以 \(\mathrm{J}(\rho). \varphi|\mathrm{T}\) 在 \(\mathrm{W}\) 下反不变。此外，根据《谱理论》和定理 1，族 \((\chi_{\lambda})_{\lambda \in \mathrm{X}_{++}}\) 是 \(\mathbf{G}\) 上中心代表函数空间的代数基，也是 \(\mathrm{ZL}^2(\mathrm{G})\)（即 \(\mathbf{G}\) 上平方可积中心函数的等价类空间）的正交归一基。

因此，由命题 3 和定理 2 得到：

推论 3. 将 G 上任意连续中心函数 \(\varphi\) 映为函数 \(w(\mathrm{G})^{-1/2}\mathrm{J}(\rho)(\varphi|\mathrm{T})\) 的映射，诱导出从 G 上中心代表函数空间到 \(\Theta(\mathrm{T})\) 的反不变元空间的同构；它由连续性延拓为从 \(\mathrm{ZL}^2(\mathrm{G})\) 到 \(\mathrm{L}^2(\mathrm{T})\) 的反不变元子空间的 Hilbert 空间同构。

推论 4. 令 \(\varphi\) 为 G 上的连续中心函数。则

\[
\int_ {\mathrm{G}} \overline {{\chi_ {\lambda} (g)}} \varphi (g) d g = \int_ {\mathrm{T}} \overline {{\lambda (t)}} \prod_ {\alpha > 0} (1 - \alpha (t) ^ {- 1}) \varphi (t) d t = \int_ {\mathrm{T}} \overline {{\lambda \rho (t)}}. \varphi (t) \mathrm{J} (\rho) (t) d t.
\]

事实上，由引理 4 和定理 2，

\[
\begin{array}{c} \int_ {\mathrm{G}} \overline {{\chi_ {\lambda} (g)}} \varphi (g) d g = \frac {1}{w (\mathrm{G})} \int_ {\mathrm{T}} \overline {{\chi_ {\lambda} (t) \mathrm{J} (\rho) (t)}} (\varphi (t) \mathrm{J} (\rho) (t)) d t \\ = \frac {1}{w (\mathrm{G})} \int_ {\mathrm{T}} \overline {{\mathrm{J} (\lambda \rho) (t)}} \varphi (t) \mathrm{J} (\rho) (t) d t. \end{array}
\]

但是函数 \(t \mapsto \varphi(t)J(\rho)(t)\) 是反不变的，而 \(\frac{1}{w(G)}J(\lambda\rho)\) 是将 \(\lambda\rho\) 投影到 \(L^{2}(T)\) 的反不变元子空间上的正交投影，因此

\[
\frac {1}{w (\mathrm{G})} \int_ {\mathrm{T}} \overline {{\mathrm{J} (\lambda \rho) (t)}} \varphi (t) \mathrm{J} (\rho) (t) d t = \int_ {\mathrm{T}} \overline {{\lambda \rho (t)}} \varphi (t) \mathrm{J} (\rho) (t) d t;
\]

最后，由公式（3），有 \(\overline{\rho(t)}\mathrm{J}(\rho)(t) = \prod_{\alpha >0}(1 - \alpha (t)^{-1})\)，故推论得证。

备注。1) 对于所有 \(w \in \mathrm{W}\)，置 \(\rho_w = {}^w\rho /\rho\)；则

\[
\sum_ {w} \varepsilon (w) \rho_ {w} = \prod_ {\alpha > 0} (1 - \alpha^ {- 1}) = \rho^ {- 2} \prod_ {\alpha > 0} (\alpha - 1).\tag{6}
\]

若 \(t\) 是 \(\mathrm{T}\) 的正则元，则由（5）得到

\[
\chi_ {\lambda} (t) = \frac {\sum_ {w} \varepsilon (w) ^ {w} \lambda (t) \rho_ {w} (t)}{\sum_ {w} \varepsilon (w) \rho_ {w} (t)} = \frac {\sum_ {w} \varepsilon (w) ^ {w} \lambda (t) \rho_ {w} (t)}{\prod_ {\alpha > 0} (1 - \alpha (t) ^ {- 1})}.\tag{7}
\]

注意，\(\rho_{w}\) 是根的整数系数线性组合，因而即使不假设 \(\rho \in \mathrm{X}(T)\)，它仍属于 X(T)。由此可知，公式（7）无需假设 \(\rho \in \mathrm{X}(T)\) 也成立：事实上，为证明这一点，只需将 G 换成适当的连通覆盖，便可归结为推论 2。

\begin{enumerate}
\def\labelenumi{\arabic{enumi})}
\setcounter{enumi}{1}
\item
  同样，推论 4 的第一个等式无需假设 \(\rho \in \mathrm{X}(\mathrm{T})\) 也成立。
\item
  可以由相应的无穷小断言（第 VIII 章，第 9 节，第 1 号，定理 1）推出定理 2；下一号的定理 3 也同样如此（它是上述引文第 2 号定理 2 的类似结果）。
\end{enumerate}

